% generated by src/report.py from docs/generated/calibration_record.json
% PROTOCOL.md §4.1. Regenerate after adding any row; do not edit this file by hand.
\begin{table}[H]
\caption{کارنامه‌ی قاعده‌ی «کالیبراسیون هم‌تراز»: هیچ تفاوتی میان دو مدل به بازنمایی نسبت
داده نمی‌شود مگر آنکه پس از هم‌تراز کردن نقاط تصمیم باقی بماند. این قاعده تا کنون
3 پیش‌بینی را رد یا تأیید کرده است و — که مهم‌تر است — در \emph{هر دو جهت}
عمل کرده: 1 مورد از آن‌ها تفاوت را 	extbf{تأیید} کرده است. بدون چنین موردی، قاعده
از استدلال جهت‌دار قابل تشخیص نمی‌بود.}
\centering
\begin{tabular}{|p{4.3cm}|p{2.6cm}|p{5.4cm}|c|}
\hline
    \textbf{ادعا} & \textbf{خوانش اولیه} & \textbf{پس از کالیبراسیون هم‌تراز} & \textbf{نتیجه} \\ \hline
    فروپاشی بازخوانی \lr{Mild} (\lr{F1}) & محدودیت ظرفیت & \textbf{کالیبراسیون} — ۵٫۴٪ به ۷۷٫۸٪ تنها با جابه‌جایی نقاط برش & رد شد \\ \hline
    برتری \lr{EfficientNet} بر \lr{DenseNet} (\lr{F3}) & بازنمایی بهتر & \textbf{جای‌گذاری مرز} — برتری از میان رفت و یک خانه تغییر علامت داد & رد شد \\ \hline
    تفاوت \lr{Moderate} میان دو منبع تصویر (\lr{E12}) & یک اثر مرزی دیگر & \textbf{تفاوت واقعی} — از ۶٫۸− به ۶٫۵− ، عملاً بدون تغییر & \textbf{تأیید شد} \\ \hline
\end{tabular}
\label{tab:calibration-record}
\end{table}
